\section{Method}
\label{sec:method}

The input is read once, as bytes. A line ends at a newline; a word is a maximal run of
bytes that are not one of the six ASCII separators (space, tab, newline, carriage
return, vertical tab and form feed); a byte count is simply the length of the input.
Carriage return is a separator rather than a line break, which is what the POSIX
specification asks for~\cite{posix2024} and what makes CRLF input agree with LF input.

\begin{table}[h]
  \centering
  \caption{Fixtures used by the test suite.}
  \label{tab:fixtures}
  \begin{tabular}{llrrr}
\toprule
Input & Description & Lines & Words & Bytes \\
\midrule
(nothing) & empty input & 0 & 0 & 0 \\
\texttt{a b} & two words, no newline & 0 & 2 & 3 \\
\texttt{a b} + newline & one terminated line & 1 & 2 & 4 \\
\texttt{h\'ello w\"orld} + newline & UTF-8 accents & 1 & 2 & 14 \\
\texttt{a b} and \texttt{c} & two terminated lines & 2 & 3 & 6 \\
\bottomrule
\end{tabular}

\end{table}

Table~\ref{tab:fixtures} shows why the byte column is the one that catches mistakes:
the accented text has the same number of characters as the ASCII text plus one, but
two more bytes.
